% !TEX program = lualatex
% LuaLaTeX でコンパイル: latexmk -lualatex presentation_revised.tex
% デザインは ../../1027_AMPL/slides/slides.tex に合わせている
\documentclass[aspectratio=169,11pt]{beamer}

\usetheme[progressbar=frametitle,numbering=fraction,sectionpage=none]{metropolis}

% 配色：爽やかな青
\definecolor{BlueMain}{HTML}{1E6BC5}    % フレームタイトル・図の線
\definecolor{BlueAccent}{HTML}{0284C7}  % 強調文字
\definecolor{BlueLight}{HTML}{38BDF8}   % プログレスバー・区切り線
\definecolor{BlueText}{HTML}{1E293B}    % 本文
\definecolor{RailGreen}{HTML}{2E8B57}
\definecolor{RoadOrange}{HTML}{F97316}
\definecolor{MidGray}{HTML}{64748B}
% 図・表で使う別名
\colorlet{TitleBlue}{BlueMain}
\colorlet{TextNavy}{BlueText}
\colorlet{LightBlue}{BlueLight!15}
\setbeamercolor{normal text}{fg=BlueText, bg=white}
\setbeamercolor{frametitle}{fg=white, bg=BlueMain}
\setbeamercolor{palette primary}{fg=white, bg=BlueMain}
\setbeamercolor{progress bar}{fg=BlueLight, bg=BlueMain!20}
\setbeamercolor{title separator}{fg=BlueLight}
\setbeamercolor{alerted text}{fg=BlueAccent}
\setbeamercolor{block title}{fg=BlueMain, bg=BlueLight!12}
\setbeamercolor{block body}{fg=BlueText, bg=BlueLight!12}
\setbeamercolor{block title alerted}{fg=BlueMain, bg=BlueMain!10}
\setbeamercolor{block body alerted}{fg=BlueText, bg=BlueMain!10}
\usepackage{luatexja-fontspec}
% 和文はゴシック。本文は Regular（Fira Sans Light に合わせる）、太字は Medium（Fira Sans Regular に合わせる）
\setmainjfont{HaranoAjiGothic-Regular.otf}[BoldFont=HaranoAjiGothic-Medium.otf]
\setsansjfont{HaranoAjiGothic-Regular.otf}[BoldFont=HaranoAjiGothic-Medium.otf]
\usepackage{amsmath,amssymb}
\usepackage{graphicx}
\usepackage{booktabs}
\usepackage{array}
\usepackage{colortbl}
\usepackage{listings}
\usepackage{tikz}
\usetikzlibrary{positioning,calc,patterns}
\usepackage{pgfplots}
\pgfplotsset{compat=1.18}
\renewcommand{\arraystretch}{1.05}

% ---------- AMPL コード ----------
\lstdefinelanguage{AMPL}{
  morekeywords={set,param,var,minimize,maximize,subject,to,s.t.,sum,in,within,binary,
    integer,default,and,Infinity},
  sensitive=true,
  morecomment=[l]{\#},
}
\lstdefinestyle{amplcode}{
  language=AMPL,
  basicstyle=\ttfamily\fontsize{7}{8.6}\selectfont,
  numbers=left,
  numberstyle=\tiny\color{MidGray},
  numbersep=4pt,
  xleftmargin=10pt,
  keywordstyle=\color{BlueAccent}\bfseries,
  commentstyle=\color{BlueMain!60!gray},
  breaklines=true,
  columns=fullflexible,
  keepspaces=true,
  showstringspaces=false,
}

% metropolis のタイトルページは高さ \paperheight の minipage を使うため Overfull \vbox が出る。高さだけ詰めて再定義する
\makeatletter
\setbeamertemplate{title page}{
  \begin{minipage}[b][\dimexpr\paperheight-16pt\relax]{\textwidth}
    \vfill%
    \ifx\inserttitle\@empty\else\usebeamertemplate*{title}\fi
    \ifx\insertsubtitle\@empty\else\usebeamertemplate*{subtitle}\fi
    \usebeamertemplate*{title separator}
    \ifx\beamer@shortauthor\@empty\else\usebeamertemplate*{author}\fi
    \ifx\insertdate\@empty\else\usebeamertemplate*{date}\fi
    \ifx\insertinstitute\@empty\else\usebeamertemplate*{institute}\fi
    \vfill
    \vspace*{1mm}
  \end{minipage}
}
\makeatother

% ---------- Helpers ----------
% 発表でさっと飛ばすページ（記号・コード・付録）：タイトル帯を薄い色にする
\setbeamerfont{block body}{size=\small}
\setbeamerfont{block body alerted}{size=\small}
\setbeamerfont{block title}{size=\normalsize}
\setbeamerfont{block title alerted}{size=\normalsize}
\newcommand{\lightframe}{\setbeamercolor{frametitle}{fg=BlueMain, bg=BlueLight!22}}
% 各ページの結論
\newcommand{\ketsuron}[1]{\par\vspace{2mm}\noindent{\setlength{\fboxsep}{4pt}\colorbox{BlueMain!10}{\parbox{\dimexpr\linewidth-8pt\relax}{\small #1}}}\par}
\newcommand{\hl}[1]{\textcolor{BlueAccent}{\textbf{#1}}}
\newcommand{\pill}[2]{\tikz[baseline=-0.5ex]\node[rounded corners=2pt,fill=#1,text=white,inner xsep=5pt,inner ysep=2pt]{\scriptsize\bfseries #2};}
\newcommand{\imgdir}{../docs/images}

\title{多目的最適化（$\varepsilon$制約法）による\\ポートからDCへの輸送計画Webアプリ}
\subtitle{コストと所要日数のトレードオフを地図で可視化}
\author{2523009 恩田梨那\quad 2523019 佐藤美桜\quad 2523023 詹媛媛}
\date{データサイエンスAI実践　2026年11月}

\begin{document}

\maketitle

% =====================================================
% Part 1: アプリ概要と発表の流れ
% =====================================================
\section{導入}



\begin{frame}{発表の流れ}
\begin{columns}[T,totalwidth=\textwidth]
\begin{column}{0.55\textwidth}
\begin{enumerate}
\item \textbf{アプリの画面}：何ができるか
\item \textbf{問題設定}：何を、どこへ、どう運ぶか
\item \textbf{最適化モデル}：線形計画としての定式化
\item \textbf{$\varepsilon$制約法}：2つの目的を同時に扱う方法
\item \textbf{結果}：パレートフロントと3つの代表解
\item \textbf{アプリの構成と使い方}：しくみと操作
\item \textbf{まとめ}：選び方の指針と限界
\end{enumerate}
\end{column}
\begin{column}{0.40\textwidth}
\begin{block}{付録}
\footnotesize
\begin{itemize}
\item 主目的の切替（日数 vs コスト）
\item 単位・パラメータ
\item 用語集
\item AMPLコード
\end{itemize}
\end{block}
\end{column}
\end{columns}
\end{frame}

\begin{frame}{アプリの画面と機能}
\begin{columns}[T,totalwidth=\textwidth]
\begin{column}{0.62\textwidth}
\centering
% 番号は画像上の座標（横2000×縦1286に換算、左上が原点）で各パネルの見出しに合わせて置く
\begin{tikzpicture}[
  numbadge/.style={circle,fill=TitleBlue,text=white,draw=white,line width=1pt,
               inner sep=1.5pt,minimum size=5mm,font=\footnotesize\bfseries}]
\node[inner sep=0pt,anchor=south west] (screen) at (0,0) {\includegraphics[width=\linewidth]{\imgdir/app_overview.png}};
\begin{scope}[x={($(screen.south east)-(screen.south west)$)},
              y={($(screen.north west)-(screen.south west)$)}]
\node[numbadge] at (30/2000,1-172/1286) {1};
\node[numbadge] at (30/2000,1-588/1286) {2};
\node[numbadge] at (500/2000,1-172/1286) {3};
\node[numbadge] at (500/2000,1-353/1286) {4};
\node[numbadge] at (500/2000,1-985/1286) {4};
\node[numbadge] at (1240/2000,0.035) {5};
\node[anchor=west,font=\tiny\bfseries,text=TitleBlue,fill=white,inner sep=1pt]
  at (1270/2000,0.035) {$\downarrow$ 画面の下：地図・輸送一覧};
\end{scope}
\end{tikzpicture}
\end{column}
\begin{column}{0.36\textwidth}
\scriptsize
\renewcommand{\arraystretch}{1.0}
\begin{tabular}{>{\bfseries}p{0.06\linewidth}p{0.86\linewidth}}
\toprule
No. & 画面の説明\\
\midrule
1 & \hl{$\varepsilon$制約法の設定}\newline 主目的、点数、計算、初期値に戻す\\
2 & \hl{パラメータ・需要・供給}\newline コスト・速度・CO\textsubscript{2}係数や需要量を変更。エラーは赤表示\\
3 & \hl{ユートピア点・ナディア点}\newline 各目的の最良値とパレート上の最悪値\\
4 & \hl{パレートフロント・$\varepsilon$選択}\newline 点の色はCO\textsubscript{2}。スライダーで解を切り替え\\
5 & \hl{地図・輸送一覧}\newline 港からDCへの道路・鉄道フローを表示\\
\bottomrule
\end{tabular}
% 入力エラー時の動き（赤表示・再計算が必要）は「アプリの使い方」のページで説明する
\end{column}
\end{columns}
\end{frame}

% =====================================================
% Part 2: 問題設定とモデル
% =====================================================
\section{問題設定とモデル}

\begin{frame}{問題設定：4港から30DCへ、2つの輸送モード}
\begin{columns}[T,totalwidth=\textwidth]
\begin{column}{0.50\textwidth}
\begin{block}{4港から30DCへ合計15,000トン}
\begin{tabular}{lr}
Los Angeles & 3,500 t\\
Tacoma & 4,500 t\\
Houston & 5,000 t\\
Oakland & 2,000 t\\
\midrule
\textbf{合計} & \textbf{15,000 t}
\end{tabular}
\end{block}
\footnotesize 各DCの需要量の合計も15,000トン。\par\vspace{1mm}\hl{題材}：MIT SCM.275x 演習09「多目的最適化（$\varepsilon$制約法）」（後半の問題）
\end{column}
\begin{column}{0.47\textwidth}
\begin{block}{道路のみ}
\pill{RoadOrange}{ROAD}\quad \textbf{速い・高コスト・CO\textsubscript{2}が多い}
\end{block}
\vspace{1mm}
\begin{block}{鉄道＋道路}
\pill{RailGreen}{RAIL}\quad \textbf{安い・CO\textsubscript{2}が少ない・遅い}\\
\footnotesize ターミナルで1.5日待つ
\end{block}
\end{column}
\end{columns}
\vfill
\ketsuron{\hl{鉄道ターミナルの簡略化}\quad 各DCは測地線距離（地図上の最短距離）で最も近い1か所の鉄道ターミナルを使う（12か所から固定）。}
\end{frame}

{\lightframe
\begin{frame}{最適化モデル①：記号とデータ}
\begin{columns}[T,totalwidth=\textwidth]
\begin{column}{0.48\textwidth}
\begin{block}{集合}
\footnotesize
\begin{tabular}{>{\bfseries}l l}
$P$ & 港の集合（4港）\\
$D$ & 配送センターの集合（30箇所）\\
$T$ & 輸送モード（road-only / rail-road）
\end{tabular}
\end{block}
\begin{block}{決定変数}
\footnotesize
$x_{pdt}\ge0$：港 $p$ からDC $d$ へモード $t$ で運ぶ量（トン）。連続変数なので分割輸送も可能。
\end{block}
\end{column}
\begin{column}{0.48\textwidth}
\begin{block}{パラメータ}
\footnotesize
\begin{tabular}{>{\bfseries}l l}
$q_d$ & DC $d$ の需要量（トン）\\
$s_p$ & 港 $p$ の供給量（トン）\\
$Q$ & 総需要 $Q=\sum_{d}q_d=15{,}000$\\
$c_{pdt}$ & 1トンあたりの運賃（\$/t）\\
$\tau_{pdt}$ & 1トンあたりの所要日数\\
$e_{pdt}$ & 1トンあたりのCO\textsubscript{2}排出量
\end{tabular}
\end{block}
\end{column}
\end{columns}
\ketsuron{\hl{目的と評価指標}\quad 最適化するのは総輸送コスト（運賃 $c_{pdt}$ の合計）と平均所要日数（$\tau_{pdt}$）。 $e_{pdt}$ は目的関数に入れず、選ばれた輸送量から結果としてCO\textsubscript{2}を集計する。}
\end{frame}
}

\begin{frame}{最適化モデル②：目的関数と制約}
\begin{columns}[T,totalwidth=\textwidth]
\begin{column}{0.48\textwidth}
\begin{block}{2つの目的関数（どちらも最小化）}
\vspace{-2mm}
\begin{align*}
\text{TotalCost}(x)&=\textstyle\sum_{p,d,t}c_{pdt}\,x_{pdt}\\
\text{AvgDuration}(x)&=\textstyle\sum_{p,d,t}\frac{\tau_{pdt}}{Q}\,x_{pdt}
\end{align*}
\end{block}
\begin{block}{目的の定義}
\scriptsize
\hl{総輸送コスト}：運賃の合計。$c_{pdt}$ は区間ごとに $\max$(単価×距離, 最低料金)、鉄道＋道路は2区間の和。積み替え・保管・在庫・時間価値は含まない。\par
\hl{平均所要日数}：各経路の日数を輸送量で重み付けして総需要 $Q$ で割った\textbf{輸送量加重平均}。
\end{block}
\end{column}
\begin{column}{0.48\textwidth}
\begin{block}{制約}
\vspace{-2mm}
\begin{align*}
\textstyle\sum_{p,t}x_{pdt}&\ge q_d\quad(\forall d)\\
\textstyle\sum_{d,t}x_{pdt}&\le s_p\quad(\forall p)\\
x_{pdt}&\ge0
\end{align*}
\footnotesize
上から、各DCの需要を満たす／各港の供給量を超えて出荷しない／負の輸送量は認めない。
\end{block}
\end{column}
\end{columns}
\ketsuron{目的関数・制約・変数がすべて線形 $\Rightarrow$ \hl{線形計画問題}（高速に厳密解が得られる）}
\end{frame}

% =====================================================
% Part 3: $\varepsilon$制約法
% =====================================================
\section{$\varepsilon$制約法}

\begin{frame}{パレート最適とは：どちらも悪くならずに良くできない解}
\begin{columns}[T,totalwidth=\textwidth]
\begin{column}{0.50\textwidth}
\centering
\begin{tikzpicture}[scale=0.85,>=stealth,font=\scriptsize]
  \def\front{(0.8,3.6) .. controls (1.2,1.5) and (2.8,0.95) .. (4.7,0.85)}
  % 実現できる解の範囲
  \fill[BlueLight!15] \front -- (4.9,3.9) -- (0.8,3.9) -- cycle;
  % A（フロント上の点）
  \path \front node[pos=0.33,inner sep=0pt] (A) {};
  % A に支配される領域
  \fill[RoadOrange!18] (A) rectangle (4.9,3.9);
  \draw[RoadOrange,dashed] (A |- 0,3.9) -- (A) -- (4.9,0 |- A);
  \node[RoadOrange] at (3.9,3.55) {A に支配される範囲};
  \node[MidGray,align=center] at (1.45,3.45) {実現できる解};
  % 軸
  \draw[->] (0,0) -- (5.2,0) node[below left] {コスト};
  \draw[->] (0,0) -- (0,4.2) node[left] {日数};
  % フロント
  \draw[BlueMain,line width=1.5pt] \front;
  \node[BlueMain,above] at (4.3,0.95) {パレートフロント};
  % 点
  \fill[BlueMain] (A) circle (2.6pt);
  \node[BlueMain,below left] at (A) {A};
  \fill[MidGray] (3.2,2.7) circle (2.4pt) node[right,text=MidGray] {B（A より高くて遅い）};
  \fill[RailGreen] (0.8,0.85) circle (2.4pt) node[below right,text=RailGreen] {ユートピア点（実現できない）};
\end{tikzpicture}
\end{column}
\begin{column}{0.47\textwidth}
\begin{block}{定義（2目的とも最小化）}
\footnotesize
\begin{itemize}\setlength{\itemsep}{1pt}
\item \hl{支配}：A が B より\textbf{コストも日数も同じか良く}、どちらかは\textbf{真に良い}
\item \hl{パレート最適解}：どの解にも支配されない解（図の A）
\item \hl{パレートフロント}：パレート最適解を目的値で並べた曲線
\end{itemize}
\end{block}
{\footnotesize\hl{$\varepsilon$制約法}：コストに上限を置いて日数を最小化した解はフロント上の点（最適解が1つならパレート最適）。上限を動かすとフロントをなぞれる。\par}
\end{column}
\end{columns}
\ketsuron{フロント上では\hl{一方を良くすると他方が必ず悪くなる}。どの点を選ぶかは\hl{意思決定}で決める}
\end{frame}

\begin{frame}{$\varepsilon$制約法のしくみ：片方を目的に、片方を上限に}
\begin{columns}[T,totalwidth=\textwidth]
\begin{column}{0.46\textwidth}
\centering
\begin{tikzpicture}[scale=0.5,>=stealth,font=\scriptsize]
  \coordinate (origin) at (0,0);
  \coordinate (top) at (0,4.0);
  % axes
  \draw[->] (0,0) -- (4.9,0) node[right] {コスト};
  \draw[->] (0,0) -- (0,4.3) node[left] {日数};
  % Pareto front
  \draw[TitleBlue,line width=1.3pt]
    (0.8,3.6) .. controls (1.2,1.6) and (2.6,1.0) .. (4.2,0.9)
    node[pos=0.5,inner sep=0pt] (m) {};
  % utopia / nadir
  \node[circle,fill=RailGreen,inner sep=2pt,label={[text=RailGreen]below right:ユートピア}] at (0.8,0.9) {};
  \node[circle,fill=MidGray,inner sep=2pt,label={[text=MidGray]right:ナディア}] at (4.2,3.6) {};
  % cost cap
  \draw[RoadOrange,thick,dashed] (m |- origin) -- (m |- top)
    node[above,text=RoadOrange,align=center] {コスト上限\\$C^U+\varepsilon(C^N-C^U)$};
  \fill[RoadOrange] (m) circle (2.2pt);
  \node[RoadOrange,above right=1pt of m] {この上限での最適解};
  % ticks
  \draw (0.8,0.06) -- (0.8,-0.06) node[below] {$C^U$};
  \draw (4.2,0.06) -- (4.2,-0.06) node[below] {$C^N$};
  \draw (0.06,0.9) -- (-0.06,0.9) node[left] {$D^U$};
  \draw (0.06,3.6) -- (-0.06,3.6) node[left] {$D^N$};
  \node[below=9pt] at (0.8,-0.1) {\tiny $\varepsilon=0$};
  \node[below=9pt] at (4.2,-0.1) {\tiny $\varepsilon=1$};
\end{tikzpicture}
\begin{block}{定式化（主目的＝平均所要日数）}
\scriptsize
\setlength{\abovedisplayskip}{2pt}\setlength{\belowdisplayskip}{2pt}
\vspace{-2mm}
\begin{align*}
\min\ & \text{AvgDuration}(x)\\
\text{s.t.}\ & \text{TotalCost}(x)\le C^U+\varepsilon(C^N-C^U)
\end{align*}
さらに需要・供給制約を加える。$\varepsilon=0$ でコスト最小、$\varepsilon=1$ で日数最短。
\end{block}
\end{column}
\begin{column}{0.50\textwidth}
\begin{block}{考え方}
\footnotesize
主目的を1つだけ最小化し、もう片方は「上限 $\varepsilon$ 以下」という制約にする。上限を少しずつ緩めて繰り返し解くと、トレードオフ曲線が描ける。
\end{block}
\begin{enumerate}\footnotesize
\setlength{\itemsep}{0pt}\setlength{\topsep}{2pt}
\item 各目的の最良値（ユートピア点）と、もう片方の値（ナディア点）を求める
\item 値を $[0,1]$ に正規化
\item 主目的（日数）を最小化
\item コストを上限 $C^U+\varepsilon(C^N-C^U)$ で制約
\item $\varepsilon$ を0から1へ動かして繰り返す
\end{enumerate}
\end{column}
\end{columns}
\end{frame}

\begin{frame}{ユートピア点とナディア点：$\varepsilon$ の範囲を決める}
% 画面画像は横長（約9:1）なので、段組みの外に横幅いっぱいで置く
\centering
\includegraphics[width=\textwidth]{\imgdir/utopia_nadir.png}\par
\vspace{-1mm}
\begin{columns}[T,totalwidth=0.97\textwidth]
\begin{column}{0.50\textwidth}
\centering\footnotesize
\renewcommand{\arraystretch}{0.95}
\begin{tabular}{lr}
\toprule
\rowcolor{LightBlue}
\textbf{区分}&\textbf{値}\\
\midrule
ユートピア：最小コスト $C^U$ & \$1,311,311\\
ユートピア：最短日数 $D^U$ & 2.50日\\
\midrule
ナディア：最大コスト $C^N$ & \$1,812,413\\
ナディア：最長日数 $D^N$ & 5.32日\\
\bottomrule
\end{tabular}
\vspace{1mm}
\begin{alertblock}{ユートピア点は実現できない}
\scriptsize
コスト最小の解は5.32日かかり、日数最短の解は\$1,812,413かかる。両方の最良値を同時には達成できない。
\end{alertblock}
\end{column}
\begin{column}{0.46\textwidth}
\begin{alertblock}{範囲の大きさ}
\scriptsize
コスト差 \$501,102（約38.2\%）／日数差 2.82日（約53\%）
\end{alertblock}
\begin{block}{定義と正規化}
\scriptsize
\setlength{\abovedisplayskip}{1pt}\setlength{\belowdisplayskip}{0pt}
\hl{ユートピア点}（理想点）：各目的の最良値の組\par
\hl{ナディア点}：\textbf{パレート最適解の中で}の最悪値の組\par
\vspace{-1mm}
\[
\bar f(x)=\frac{f(x)-f^U}{f^N-f^U}\in[0,1]
\]
\end{block}
\end{column}
\end{columns}
\end{frame}

% =====================================================
% Part 4: 結果
% =====================================================
\section{結果}

\begin{frame}{結果：$\varepsilon$ ごとのコスト・日数・CO\textsubscript{2}（主目的＝日数）}
\footnotesize
\renewcommand{\arraystretch}{0.82}
\vspace*{-2mm}
\begin{center}
\begin{tabular}{c r r r c c}
\toprule
$\varepsilon$ & コスト & 平均日数 & CO\textsubscript{2} & 道路 & 鉄道＋道路\\
\midrule
0.0 & \$1,311,312 & 5.32日 & 1,151t & 34.1\% & 65.9\%\\
\rowcolor{LightBlue}0.1 & \$1,361,422 & 4.49日 & 1,556t & 55.5\% & 44.5\%\\
\rowcolor{LightBlue}0.2 & \$1,411,532 & 4.07日 & 1,804t & 67.4\% & 32.6\%\\
\rowcolor{LightBlue}0.3 & \$1,461,642 & 3.79日 & 2,001t & 75.7\% & 24.3\%\\
0.4 & \$1,511,752 & 3.57日 & 2,197t & 80.0\% & 20.0\%\\
0.5 & \$1,561,862 & 3.38日 & 2,386t & 83.7\% & 16.3\%\\
0.6 & \$1,611,973 & 3.19日 & 2,575t & 87.1\% & 12.9\%\\
0.7 & \$1,662,083 & 3.00日 & 2,763t & 90.5\% & 9.5\%\\
0.8 & \$1,712,193 & 2.82日 & 2,940t & 94.0\% & 6.0\%\\
0.9 & \$1,762,303 & 2.66日 & 3,107t & 97.1\% & 2.9\%\\
1.0 & \$1,812,413 & 2.50日 & 3,272t & 100.0\% & 0.0\%\\
\bottomrule
\end{tabular}
\par{\scriptsize\color{BlueText!65} ※ $\varepsilon=0$ のコストが $C^U$（\$1,311,311）と\$1違うのは、数値誤差対策でコスト上限を相対 $10^{-7}$ 緩めて解くため\par}
\end{center}
\vspace{-3mm}
\ketsuron{\hl{読み方}\quad $\varepsilon=0$ は安いが遅い端、$\varepsilon=1$ は速いが高い端。水色の行（$\varepsilon=0.1\text{〜}0.3$）は日数が最も大きく縮む区間。}
\end{frame}

\begin{frame}{パレートフロント：読み方のポイント}
\begin{columns}[T,totalwidth=\textwidth]
\begin{column}{0.50\textwidth}
\includegraphics[width=\linewidth,height=0.58\textheight,keepaspectratio]{\imgdir/pareto_front.png}
\ketsuron{\hl{収穫逓減}\quad $\varepsilon$ が0.1増えるたびコストは約+\$50,110。 1日短縮の追加コストは\hl{約\$6万/日 $\to$ 約\$32万/日}（5倍超）。}
\end{column}
\begin{column}{0.47\textwidth}
% 箇条書きの中はテンプレートで \small になるため、\footnotesize は enumerate の中に書く
\begin{enumerate}\scriptsize
\setlength{\itemsep}{1pt}\setlength{\topsep}{0pt}
\item \hl{曲線は右下がり}\\
左上は安いが遅く、右下は速いが高い。曲線上の点はどれも最適で、方針によって選ぶ。曲線より左下は実現できない。
\item \hl{最初の追加予算ほど効く}\\
追加コスト全体（\$501,102）の30\%を使う $\varepsilon=0.3$ で、日数の短縮幅2.82日のうち約54\%（1.53日）を獲得。最後の区間（0.9$\to$1.0）はわずか0.16日。
\item \hl{速くするほどCO\textsubscript{2}が増える}\\
鉄道割合は65.9\%$\to$0\%、CO\textsubscript{2}は1,151t$\to$3,272t（約2.8倍）。点の濃さがCO\textsubscript{2}を示す。
\item \hl{鉄道が最後まで残るのはTacoma発}\\
$\varepsilon=0.5$ ではTacoma港からの鉄道輸送2,451tだけが残り、他の港は道路のみ。
\end{enumerate}
\end{column}
\end{columns}
\end{frame}

\begin{frame}{なぜ収穫逓減になるのか：フロントは凸の折れ線}
\begin{columns}[T,totalwidth=\textwidth]
\begin{column}{0.52\textwidth}
\centering
\begin{tikzpicture}
\begin{axis}[
  width=\linewidth, height=4.6cm,
  ybar, bar width=7pt,
  ymin=0, ymax=36,
  xlabel={$\varepsilon$ の区間}, ylabel={万ドル／日},
  xtick=data,
  xticklabels={0--.1,,.2--.3,,.4--.5,,.6--.7,,.8--.9,},
  tick label style={font=\scriptsize}, label style={font=\scriptsize},
  nodes near coords, nodes near coords style={font=\tiny},
  point meta=explicit symbolic,
  axis lines*=left, ymajorgrids, grid style={BlueLight!30},
  enlarge x limits=0.06,
]
\addplot[fill=BlueMain!70,draw=none] coordinates {
  (1,6.0)[6.0] (2,12.0)[12.0] (3,17.7)[17.7] (4,23.4)[23.4] (5,25.4)[25.4]
  (6,26.4)[26.4] (7,27.0)[27.0] (8,28.5)[28.5] (9,30.8)[30.8] (10,31.9)[31.9]};
\end{axis}
\end{tikzpicture}\par
{\scriptsize 1日短縮あたりの追加コスト（各区間でコスト $+\$50{,}110$）\par}
\end{column}
\begin{column}{0.45\textwidth}
\begin{block}{線形計画の性質}
\footnotesize
コスト上限 $C$ での最短日数 $D(C)$ は、LP では\hl{凸で区分線形}（折れ線）。
\end{block}
\begin{block}{意味}
\footnotesize
\begin{itemize}\setlength{\itemsep}{1pt}
\item 傾き $-D'(C)$（1ドルで減る日数）は、凸なので\hl{予算を足すほど小さい}
\item その逆数（1日短縮の追加コスト）は\hl{必ず増える}（左図）
\item 傾き＝コスト上限の\hl{双対変数}
\end{itemize}
\end{block}
\end{column}
\end{columns}
\ketsuron{収穫逓減は偶然ではなく、\hl{LPのパレートフロントが凸であることから必ず起きる}}
\end{frame}

\begin{frame}{代表解の比較：$\varepsilon=0$、0.5、1}
\begin{columns}[T,totalwidth=\textwidth]
\begin{column}{0.31\textwidth}
\centering
\textbf{$\varepsilon=0$（コスト最小）}\\
\includegraphics[width=\linewidth,height=0.46\textheight,keepaspectratio]{\imgdir/epsilon_result_0.png}
\vspace{2mm}
\footnotesize
\begin{tabular}{lr}
コスト & \$1,311,312\\
平均日数 & 5.32日\\
CO\textsubscript{2} & 1,151t\\
道路：鉄道 & 34.1：65.9\%
\end{tabular}
\end{column}
\begin{column}{0.31\textwidth}
\centering
\textbf{$\varepsilon=0.5$（中間）}\\
\includegraphics[width=\linewidth,height=0.46\textheight,keepaspectratio]{\imgdir/epsilon_result_0.5.png}
\vspace{2mm}
\footnotesize
\begin{tabular}{lr}
コスト & \$1,561,862\\
平均日数 & 3.38日\\
CO\textsubscript{2} & 2,386t\\
道路：鉄道 & 83.7：16.3\%
\end{tabular}
\end{column}
\begin{column}{0.31\textwidth}
\centering
\textbf{$\varepsilon=1$（日数最短）}\\
\includegraphics[width=\linewidth,height=0.46\textheight,keepaspectratio]{\imgdir/epsilon_result_1.png}
\vspace{2mm}
\footnotesize
\begin{tabular}{lr}
コスト & \$1,812,413\\
平均日数 & 2.50日\\
CO\textsubscript{2} & 3,272t\\
道路：鉄道 & 100：0\%
\end{tabular}
\end{column}
\end{columns}
\vfill
\ketsuron{\hl{注}\quad 主目的は平均所要日数。同じ $\varepsilon=0.5$ でも総輸送コストを主目的にすると別の点（\$1,435,456・3.91日）になる（付録参照）。}
\end{frame}

\begin{frame}{地図で見る両端の輸送計画}
\begin{columns}[T,totalwidth=\textwidth]
\begin{column}{0.48\textwidth}
\centering
\pill{TitleBlue}{$\varepsilon=0$ コスト最小}\\[2mm]
\includegraphics[width=\linewidth,height=0.55\textheight,keepaspectratio]{\imgdir/map_eps0.png}\\
\small 緑（鉄道）が多い／CO\textsubscript{2} 1,151t
\end{column}
\begin{column}{0.48\textwidth}
\centering
\pill{TitleBlue}{$\varepsilon=1$ 日数最短}\\[2mm]
\includegraphics[width=\linewidth,height=0.55\textheight,keepaspectratio]{\imgdir/map_eps1.png}\\
\small すべてオレンジ（道路）／CO\textsubscript{2} 3,272t
\end{column}
\end{columns}
\vfill
\ketsuron{\hl{環境への代償}\quad 所要日数を短縮するほど鉄道割合が減り、CO\textsubscript{2}は約\hl{2.8倍}になる。}
\end{frame}

\begin{frame}{意思決定の指針：どの $\varepsilon$ を選ぶか}
\footnotesize
\begin{center}
\begin{tabular}{>{\bfseries}l c r r r}
\toprule
\rowcolor{LightBlue}
\textbf{重視すること} & 目安の$\varepsilon$ & コスト & 平均日数 & CO\textsubscript{2}\\
\midrule
コスト・環境 & 0〜0.1 & \$1.31M〜1.36M & 5.32〜4.49日 & 1,151〜1,556t\\
バランス & 0.3 & \$1.46M & 3.79日 & 2,001t\\
納期優先 & 0.7〜1.0 & \$1.66M〜1.81M & 3.00〜2.50日 & 2,763〜3,272t\\
\bottomrule
\end{tabular}
\end{center}
\vspace{2mm}
\begin{columns}[T,totalwidth=\textwidth]
\begin{column}{0.48\textwidth}
\begin{block}{1日短縮あたりの追加コスト}
\footnotesize
\begin{itemize}
\item $\varepsilon=0\to0.3$：+\$150,330で1.53日短縮（約\$10万/日）
\item $\varepsilon=0.3\to1.0$：+\$350,771で1.29日短縮（約\$27万/日）
\end{itemize}
\end{block}
\end{column}
\begin{column}{0.48\textwidth}
\ketsuron{\hl{使い方}\quad 「1日の短縮に\$10万を払えるか」という基準があれば、その金額に合う点がそのまま選べる。目安は表示値から求めたもので、最終判断は業務上の納期価値による。}
\end{column}
\end{columns}
\end{frame}

% =====================================================
\begin{frame}{結論：安く運ぶと、速くは運べない}
\begin{columns}[T,totalwidth=\textwidth]
\begin{column}{0.31\textwidth}
\begin{block}{総輸送コスト}
\centering
{\Large\bfseries +38\%}\\[1mm]
\footnotesize \$1,311,311 $\to$ \$1,812,413\\（差額 \$501,102）
\end{block}
\end{column}
\begin{column}{0.31\textwidth}
\begin{block}{平均所要日数}
\centering
{\Large\bfseries 5.32日 $\to$ 2.50日}\\[1mm]
\footnotesize 約53\%短縮（コスト最小解 $\to$ 日数最短解）
\end{block}
\end{column}
\begin{column}{0.31\textwidth}
\begin{block}{CO\textsubscript{2}排出量}
\centering
{\Large\bfseries 約2.8倍}\\[1mm]
\footnotesize 1,151t $\to$ 3,272t\\（鉄道が減るため）
\end{block}
\end{column}
\end{columns}
\vspace{3mm}
\begin{itemize}
\item 4港から30配送センター（DC）へ\textbf{15,000トン}を運ぶ計画を対象にした
\item \textbf{$\varepsilon$制約法}で、コストと日数のトレードオフ曲線（パレートフロント）を作成
\item 追加コストの\textbf{最初の30\%で、短縮幅の約54\%}を得られる（収穫逓減）
\end{itemize}
\ketsuron{\hl{このアプリの役割}\quad 「単一の最適解」ではなく、\hl{どこまで払ってどこまで速くするか}を選ぶための地図を示す。}
\end{frame}

% Part 5: アプリ
% =====================================================
\section{アプリの使い方}



\begin{frame}{アプリの構成：画面・API・最適化の3層}
\centering
\begin{tikzpicture}[
  box/.style={draw=TitleBlue,fill=LightBlue!60,rounded corners=3pt,align=center,
              text width=33mm,minimum height=17mm,font=\scriptsize},
  pre/.style={draw=MidGray,dashed,fill=white,rounded corners=3pt,align=center,
              text width=36mm,minimum height=10mm,font=\scriptsize},
  arr/.style={->,>=stealth,thick,TitleBlue},
  lab/.style={font=\tiny,text=TextNavy,align=center}]
\node[box,text width=29mm] (ui) {\textbf{画面（React）}\\[0.5mm] グラフ：Recharts\\ 地図：Leaflet＋OpenStreetMap};
\node[box,right=11mm of ui,text width=42mm] (api) {\textbf{API（FastAPI）}\\[0.5mm] /api/network（初期データ）\\ /api/frontier（パレートフロント）\\ /api/solve（指定した $\varepsilon$）};
\node[box,right=12mm of api,text width=29mm] (opt) {\textbf{最適化（AMPL＋HiGHS）}\\[0.5mm] amplpy でデータを渡し\\ 線形計画を解く};
\draw[arr] ([yshift=2mm]ui.east) -- node[lab,above]{入力値} ([yshift=2mm]api.west);
\draw[arr] ([yshift=-2mm]api.west) -- node[lab,below]{結果（JSON）} ([yshift=-2mm]ui.east);
\draw[arr] ([yshift=2mm]api.east) -- node[lab,above]{モデル＋データ} ([yshift=2mm]opt.west);
\draw[arr] ([yshift=-2mm]opt.east) -- node[lab,below]{解} ([yshift=-2mm]api.east);
\node[pre,below=7mm of api] (net) {\textbf{network.json}\\ 港・DC間の道路・鉄道の距離};
\node[pre,left=15mm of net] (sc) {\textbf{事前計算（1回だけ）}\\ scgraph で最短経路を計算};
\draw[arr,MidGray] (sc) -- (net);
\draw[arr,MidGray] (net) -- node[lab,right]{起動時に読み込む} (api);
\end{tikzpicture}
\vspace{2mm}
\ketsuron{\hl{スライダーが待ち時間なしで切り替わる理由}\quad 「計算」のとき、表示する点に加えて $\varepsilon$ を\hl{0.01刻みにした101点}も、AMPLの集合 $K$ を使って\hl{1回のソルバー呼び出しでまとめて解いておく}。 スライダーは計算済みの解を表示するだけなので、待ち時間なしで切り替わる。0.333 など0.01刻みでない値だけは、その1点を追加で解く（約5秒）。}
\end{frame}

\begin{frame}{使用したモデル・ソルバーと計算時間}
\small
\begin{tabular}{@{}l l >{\raggedright\arraybackslash}p{7.4cm}@{}}
\toprule
役割 & 名前 & 内容\\
\midrule
\hl{モデル} & 線形計画（LP） & 1問：変数240（4港×30DC×2モード）、制約35\newline（需要30・供給4・$\varepsilon$制約1）\\
\hl{モデリング言語} & AMPL（amplpy） & 元教材（gurobipy／MyPuLP）から書き換え\\
\hl{ソルバー} & HiGHS & AMPL に同梱のオープンソースLPソルバー\\
\bottomrule
\end{tabular}
\vspace{1mm}

\begin{tabular}{@{}l l r@{}}
\toprule
操作 & ソルバー呼び出し & 時間（作成PC）\\
\midrule
「計算」ボタン & 3回：ユートピア点 → ナディア点 → 101問まとめて & 12〜17秒\\
0.01刻みでない$\varepsilon$ & 1回 & 約5秒\\
\bottomrule
\end{tabular}
\ketsuron{LP自体は101問でも約0.1秒（scipy の HiGHS で計測した参考値）。時間の大半は\hl{AMPLの起動（1回2〜4秒）}}
\end{frame}

\begin{frame}{アプリの使い方}
\begin{columns}[T,totalwidth=\textwidth]
\begin{column}{0.38\textwidth}
\includegraphics[width=\linewidth]{\imgdir/settings_panel.png}
\begin{block}{実用上のポイント}
入力エラーは欄が赤表示になる。入力を変えたら再計算が必要。
\end{block}
\end{column}
\begin{column}{0.59\textwidth}
\begin{enumerate}\small
\setlength{\itemsep}{1pt}\setlength{\topsep}{0pt}
\item アプリを開く（起動時は初期値で計算：10〜20秒）
\item 主目的を選ぶ（初期値は所要日数の最小化）
\item 必要ならパラメータを変更して「計算」
\item パレートフロント上の点にマウスを乗せて値を確認し、クリックして選択
\item $\varepsilon$ スライダーで解を切り替える（0.01刻みの101点は計算済みなので即座に切り替わる。0.333などの細かい値は数値欄から入力し、追加で約5秒）
\item 地図と一覧表で、港・DC・輸送モード・輸送量を確認
\end{enumerate}
\end{column}
\end{columns}
\end{frame}

% =====================================================
% Part 6: まとめ
% =====================================================
\section{まとめ}

\begin{frame}{前提・限界と今後の拡張}
\begin{columns}[T,totalwidth=\textwidth]
\begin{column}{0.48\textwidth}
\begin{block}{モデルの前提と限界}
\footnotesize
\begin{itemize}\setlength{\itemsep}{0pt}
\item 各DCの鉄道ターミナルは、測地線距離で最も近い1か所に固定
\item 輸送量は連続量（端数あり）
\item 距離は事前計算済みの道路・鉄道ネットワークを使用
\item CO\textsubscript{2}は結果として表示する指標で、目的関数には含めない
\item 計算時間は1回12〜17秒程度（作成PCでの実測）
\end{itemize}
\end{block}
\end{column}
\begin{column}{0.48\textwidth}
\begin{block}{今後の拡張案}
\footnotesize
\begin{itemize}\setlength{\itemsep}{0pt}
\item CO\textsubscript{2}を第3の目的に加えた3目的への拡張
\item DCごとに複数の鉄道ターミナルから選べるようにする
\item 輸送量を整数化（トラック・コンテナ単位）
\item 需要や速度の変動に対する感度分析
\end{itemize}
\end{block}
\end{column}
\end{columns}
\ketsuron{\hl{注意}\quad CO\textsubscript{2}も同時に最小化したわけではない。日数を縮めた結果としてCO\textsubscript{2}が増える。}
\end{frame}

\begin{frame}{まとめ}
\begin{columns}[T,totalwidth=\textwidth]
\begin{column}{0.58\textwidth}
\begin{block}{わかったこと}
\begin{itemize}
\item コスト最小（5.32日）と日数最短（\$1.81M）は\hl{両立しない}
\item $\varepsilon=0.1\text{〜}0.3$ 付近が\hl{費用対効果の高い領域}（+\$15万で日数5.32$\to$3.79日）
\item 日数を縮めるほど鉄道が減り、CO\textsubscript{2}は最大\hl{約2.8倍}
\end{itemize}
\end{block}
\end{column}
\begin{column}{0.38\textwidth}
\begin{block}{アプリで実現したこと}
\footnotesize
\begin{itemize}
\item パレートフロントの可視化
\item 点の選択で地図と輸送量が連動
\item パラメータ変更と再計算
\end{itemize}
\end{block}
\end{column}
\end{columns}
\vfill
\ketsuron{意思決定を支えるのは「単一の最適解」ではなく、\hl{選択可能なパレート解}}
\end{frame}

% =====================================================
% Appendix
% =====================================================
\appendix
\section{付録}

{\lightframe
\begin{frame}{付録A：主目的を切り替えたときの $\varepsilon$ 制約}
\begin{columns}[T,totalwidth=\textwidth]
\begin{column}{0.48\textwidth}
\begin{block}{主目的：平均所要日数}
\[
\min\ \text{AvgDuration}(x)
\]
\[
\text{s.t.}\quad
\text{TotalCost}(x)\le C^U+\varepsilon(C^N-C^U)
\]
\end{block}
\footnotesize $\varepsilon=0$ はコスト最小、$\varepsilon=1$ は日数最短側の端。
\end{column}
\begin{column}{0.48\textwidth}
\begin{block}{主目的：総輸送コスト}
\[
\min\ \text{TotalCost}(x)
\]
\[
\text{s.t.}\quad
\text{AvgDuration}(x)\le D^U+\varepsilon(D^N-D^U)
\]
\end{block}
\footnotesize 同じパレートフロントを描くが、同じ $\varepsilon$ でも指す点が異なる。
\end{column}
\end{columns}
\vfill
\ketsuron{\hl{確認例：$\varepsilon=0.5$}\quad 所要日数を主目的 $\to$ \$1,561,862・3.38日\newline 総輸送コストを主目的 $\to$ \$1,435,456・3.91日}
\end{frame}
}

{\lightframe
\begin{frame}{付録B：単位とパラメータ}
\begin{columns}[T,totalwidth=\textwidth]
\begin{column}{0.48\textwidth}
\begin{block}{単位と係数}
\footnotesize
\begin{itemize}\setlength{\itemsep}{0pt}
\item 所要日数は輸送量加重平均。道路・鉄道とも1日8時間稼働
\item CO\textsubscript{2}排出量はトン
\item 速度：道路60mph、鉄道35mph
\item ターミナル滞留：1.5日
\item CO\textsubscript{2}係数：道路0.4、鉄道0.05ポンド/トン・マイル
\end{itemize}
\end{block}
\end{column}
\begin{column}{0.48\textwidth}
\begin{block}{コスト}
\footnotesize
\begin{itemize}\setlength{\itemsep}{0pt}
\item コスト単価：道路0.10、鉄道0.05（1トン・マイルあたり）
\item 最低料金：区間ごとに道路25、鉄道45（1トンあたり）
\item 通貨単位は元問題に記載がなく、ドルと仮定
\end{itemize}
\end{block}
\end{column}
\end{columns}
\end{frame}
}

{\lightframe
\begin{frame}{付録C：用語集}
\small
\begin{tabular}{>{\bfseries}p{0.22\linewidth}>{\raggedright\arraybackslash}p{0.72\linewidth}}
\toprule
用語 & 意味\\
\midrule
多目的最適化 & 複数の目的（コストと日数）を同時に扱う\\
$\varepsilon$制約法 & 主目的を最小化し、残りを「上限$\varepsilon$以下」という制約にして繰り返し解く\\
パレート最適解 & どれか1つの目的を悪くしないと、他の目的を良くできない解\\
パレートフロント & パレート最適解を目的値でグラフにしたトレードオフ曲線\\
ユートピア点 & 各目的の最良値を組み合わせた理想点。通常は実現できない\\
ナディア点 & パレート解上で各目的が最も悪くなる値を組み合わせた点\\
主目的・副目的 & $\varepsilon$制約法で最小化する目的／上限$\varepsilon$をかける目的\\
収穫逓減 & 追加量を増やしても、得られる効果がだんだん小さくなること\\
\bottomrule
\end{tabular}
\end{frame}
}

{\lightframe
\begin{frame}[fragile]{付録D：AMPLコード（モデル定義）}
% 1列では下端からはみ出すため、集合・パラメータ（左）と変数（右）の2列に分ける
\begin{columns}[T,totalwidth=\textwidth]
\begin{column}{0.44\textwidth}
\begin{lstlisting}[style=amplcode]
set K;         # 同時に解くε制約問題
set PORTS;
set DCS;
set MODES;

param demand {DCS} >= 0;
param incoming_flow {PORTS} >= 0;
param total_demand > 0;
param unit_cost {PORTS,DCS,MODES} >= 0;
param duration {PORTS,DCS,MODES} >= 0;

param w_cost {K} default 0;
param w_duration {K} default 0;
param cost_cap {K} default Infinity;
param duration_cap {K} default Infinity;
\end{lstlisting}
\end{column}
\begin{column}{0.54\textwidth}
\begin{lstlisting}[style=amplcode,firstnumber=16]

var flow {K,PORTS,DCS,MODES} >= 0;

var TotalCost {k in K} =
  sum {p in PORTS,d in DCS,t in MODES}
      unit_cost[p,d,t]*flow[k,p,d,t];

var AvgDuration {k in K} =
  sum {p in PORTS,d in DCS,t in MODES}
      duration[p,d,t]/total_demand*flow[k,p,d,t];
\end{lstlisting}
\end{column}
\end{columns}
\end{frame}
}

{\lightframe
\begin{frame}[fragile]{付録E：AMPLコード（目的・制約・$\varepsilon$制約）}
\begin{lstlisting}[style=amplcode]
minimize Objective:
  sum {k in K}
      (w_cost[k]*TotalCost[k] + w_duration[k]*AvgDuration[k]);

subject to Demand {k in K,d in DCS}:
  sum {p in PORTS,t in MODES} flow[k,p,d,t] >= demand[d];

subject to Supply {k in K,p in PORTS}:
  sum {d in DCS,t in MODES} flow[k,p,d,t] <= incoming_flow[p];

# 副目的の上限をεで段階的に設定
subject to CostCap {k in K}: TotalCost[k] <= cost_cap[k];
subject to DurationCap {k in K}:
  AvgDuration[k] <= duration_cap[k];
\end{lstlisting}
\ketsuron{\hl{所要日数を主目的}\quad $w_{\text{duration}}=1$,\ $\text{TotalCost}\le C^{U}+\varepsilon(C^{N}-C^{U})$\newline
\hl{総輸送コストを主目的}\quad $w_{\text{cost}}=1$,\ $\text{AvgDuration}\le D^{U}+\varepsilon(D^{N}-D^{U})$}
\end{frame}
}

\end{document}
